\documentclass[10pt,a4paper]{amsart}
\usepackage[margin=19mm]{geometry}
\usepackage{amsmath,amssymb,mathtools}
\usepackage[scr=rsfs,cal=euler]{mathalfa}
\usepackage{xfrac}
\usepackage[unicode,hidelinks,bookmarks=false]{hyperref}
\newcommand{\ie}{i.e.\ }
\newcommand{\cf}{cf.\ }
\newcommand{\resp}{respectively\ }
\newcommand{\Subsection}{\subsection}
\providecommand{\nobreakdash}{\nobreak\hbox{-}}
\setlength{\emergencystretch}{2em}
\multlinegap0pt
\usepackage{bm}
\usepackage{bbm}
\usepackage{dsfont}
\usepackage{xspace}
\usepackage{cancel}
\usepackage{mathtools}
\usepackage{amsthm}
\makeatletter
\@ifundefined{lemma}{\newtheorem{lemma}{Lemma}}{}
\@ifundefined{theorem}{\newtheorem{theorem}{Theorem}}{}
\makeatother
\newcommand{\ad}[1]{\mathrm{ad}(#1)}
\title[Elementary proofs of ring commutativity theorems]{Elementary proofs of ring commutativity theorems}
\author{Michael Kinyon, Desmond MacHale}
\date{Source: arXiv:2601.12599v1 (CC BY 4.0)}
\begin{document}
\maketitle

\noindent\emph{Source and licence.} Elementary proofs of ring commutativity theorems. Version arXiv:2601.12599v1 (CC BY 4.0), distributed under the Creative Commons Attribution 4.0 International licence, as recorded in the arXiv licence metadata for this exact version and shown on the abstract page https://arxiv.org/abs/2601.12599v1. Full rights notice: licenses/arxiv-2601.12599-CC-BY-4.0.txt.\par\smallskip

\noindent\textit{Editorial scope.} Counted unit: the Theorem whose source label is \texttt{Thm:exp8} in the article (source0.tex lines 501--504, source label \texttt{Thm:exp8}), delivered together with the article's own complete original proof of it (source0.tex lines 506--575, one \texttt{proof} environment of 70 lines). No mathematics is written, repaired or re-derived: statement and proof are the article's own text line for line. Required context retained uncounted: Lem:even (lines 393--395); Lem:NL (lines 428--437); Lem:equiv (lines 472--476). Every definition, display and label of those lines is retained unchanged. Omitted from this package and not redistributed: 1--392, 396--427, 438--471, 477--500, 505--505, 576--699. Nothing inside a retained excerpt is omitted or altered: every retained body is delivered exactly as the article's lines are written, so no omission receipt of any kind accompanies this package. Citations: the retained excerpts carry no citation command at all, so no bibliography entry of the article is transcribed and no reference is redistributed. Reference adaptation: none was needed and none was made; every reference inside the delivered unit resolves inside it, so no unresolved reference or citation is produced. Printed slips kept literal: no printed word, symbol, hypothesis or number of the counted statement or of its proof was altered. Layout and class: the article's own class is \texttt{bimsart} with packages including amsmath, amsthm, amssymb, bims, hyperref; the wrapper is the lane's pinned \texttt{amsart} editorial wrapper at 10pt on A4, which restates the theorem-like environments the retained text needs and adds the article's own macro definitions that the retained text needs (\texttt{\textbackslash ad}), with extra wrapper packages (bm, bbm, dsfont, xspace, cancel, mathtools). The delivered numbering is the unit's own. Assets: no retained span contains a figure environment, a graphics-inclusion command, a PDF-page inclusion, a table or a TikZ picture; no image file is copied into this package, the package declares no asset dependency and no image file is redistributed with it. Licence: version arXiv:2601.12599v1 of this article is distributed under the Creative Commons Attribution 4.0 International licence, as recorded in the arXiv licence metadata for this exact version and shown on the abstract page https://arxiv.org/abs/2601.12599v1. A full rights notice is delivered at licenses/arxiv-2601.12599-CC-BY-4.0.txt. Characters: every retained excerpt is pure ASCII, so the delivered engine, XeLaTeX with its default Latin Modern text font, reports no missing character.
\begin{lemma}\label{Lem:even}
	Let $R$ be a ring satisfying (E). Then $2\,\ad{x}=0$ for all $x\in R$.
\end{lemma}

\begin{lemma}\label{Lem:NL}
	Let $R$ be a ring satisfying (E). Then for all $x,y\in R$,
	\[
	[x,[x,y]] = 0 \implies [x,y] = 0\,,
	\]
	or equivalently,
	\[
	\ad{x}^2(y) = 0 \implies \ad{x}(y) = 0\,.
	\]
\end{lemma}

\begin{lemma}\label{Lem:equiv}
	Let $R$ be a ring, let $n$ be a positive integer and assume 
	$[x^{2^n} - x,y] = 0$ for all $x,y\in R$. Then 
	$[x^{2^{n-1}} + \cdots + x^2 + x,y] = 0$ for all $x,y\in R$.
\end{lemma}

\begin{theorem}\label{Thm:exp8}
	Let $R$ be a ring satisfying $[x^8 - x,y] = 0$ for all $x,y\in R$. Then
	$R$ is commutative.
\end{theorem}

\begin{proof}
	We will freely use Lemma \ref{Lem:even} without explicit reference.
	By Lemma \ref{Lem:equiv}, $[x^4 + x^2 + x,y] = 0$ for all $x,y\in R$, that is, $\ad{x}^4 + \ad{x}^2 + \ad{x} = 0$. 
	
	First, for all $u,v\in R$,
	\begin{align*}
		(\ad{u}^2 + \ad{v}^2)([u,v]) &= (\ad{u}^2 + \ad{[u,v]} + \ad{v}^2)([u,v]) \\
		&= (\ad{u}^2 + \ad{u}\ad{v} + \ad{v}\ad{u} + \ad{v}^2)([u,v]) \\
		&= (\ad{u} + \ad{v})^2 ([u,v]) \\
		&= \ad{u+v}^2([u,v]) \\
		&= \ad{u+v}^2([u+v,v]) \\
		&= \ad{u+v}^3(v)\,.
	\end{align*}
	Thus
	\begin{align*}
		\ad{u+v}(\ad{u}^2 + \ad{v}^2)([u,v]) &= \ad{u+v}^4(v) \\
		&= \ad{u+v}^2(v) + \ad{u+v}(v) \\
		&= \ad{u+v}([u+v,v] + v) \\
		&= \ad{u+v}(v + [u,v])\,.
	\end{align*}
	Rearranging this, we have $\ad{u+v}\ad{v}^2([u,v]) = \ad{u+v}(v + [u,v] + \ad{u}^2([u,v]))$, that is,
	\[
	\ad{u+v}\ad{v}^3(u) = \ad{u+v}(v + \ad{u}(v) + \ad{u}^3(v))\,.
	\]
	In this last equation, replace $v$ with $[u,v]$ to get
	\begin{equation}\label{Eqn:exp8-1}
		[u + [u,v], \ad{[u,v]}^3(u)] = [u + [u,v],\ad{u}(v) + \ad{u}^2(v) + \ad{u}^4{v}] = 0\,.
	\end{equation}
	
	Now let $u = [x^2,y^2 + y]$ and $v = y^2 + y$. Then 
	\begin{align*}
		[u,v] &= \ad{y^2 + y}^2(x^2) \\
		&= (\ad{y}^2+\ad{y})^2(x^2) \\
		&= (\ad{y}^4 + \ad{y}^2)(x^2) \\
		&= \ad{y}(x^2) \\
		&= [x^2,y]\,,
	\end{align*}
	and $u + [u,v] = [x^2,y^2]$. Plugging all this into \eqref{Eqn:exp8-1}, we have
	\begin{equation}\label{Eqn:exp8-2}
		[[x^2,y^2], \ad{[x^2,y]}^3([x^2,y^2 + y])] = 0\,.
	\end{equation}
	Now
	\begin{align*}
		\ad{[x^2,y]}^3([x^2,y^2 + y]) &= \ad{[x^2,y]}^2([[x^2,y],[x^2,y^2] + [x^2,y]]) \\
		&= \ad{[x^2,y]}^3([x^2,y^2]) \\
		&= \ad{[x^2,y]}^3\ad{y}^2(x^2) \\
		&= \ad{[x^2,y]}^3\ad{y}([x^2,y]) \\
		&= \ad{[x^2,y]}^3[[x^2,y],y]) \\
		&= \ad{[x^2,y]}^4(y) \\
		&= \ad{[x^2,y]}^2(y) + \ad{[x^2,y]}(y) \\
		&= \ad{[x^2,y]}\ad{y}^2(x^2) + \ad{y}^2(x^2) \\
		&= [[x^2,y],[x^2,y^2]] + [x^2,y^2]\,.
	\end{align*}
	Thus the left side of \eqref{Eqn:exp8-2} simplifies to
	\[
	[[x^2,y^2],[[x^2,y],[x^2,y^2]] + [x^2,y^2]] = [[x^2,y^2],[[x^2,y],[x^2,y^2]] = \ad{[x^2,y^2]}^2([x^2,y])\,.
	\]
	
	Therefore \eqref{Eqn:exp8-2} reduces to
	\[
	\ad{[x^2,y^2]}^2([x^2,y]) = 0
	\]
	for all $x,y\in R$. By Lemma \ref{Lem:NL},
	\[
	0 = \ad{[x^2,y^2]}([x^2,y]) = \ad{[x^2,y]}\ad{y}^2(x^2) = \ad{[x^2,y]}^2(y)\,.
	\]
	By Lemma \ref{Lem:NL} again, $\ad{[x^2,y]}(y) = 0$, 
	that is, $\ad{y}^2(x^2) = 0$. By Lemma \ref{Lem:NL} again, $\ad{y}(x^2) = [x^2,y] = \ad{x}^2(y) = 0$.
	Applying Lemma \ref{Lem:NL} one last time, $[x,y] = 0$. Therefore $R$ is commutative.
\end{proof}

\end{document}
